\section{Proof of Theorems A and B}\label{sec:deduction}

The lattice family is chosen before the volume sequence. Full convergence at each fixed volume then leaves only a volume refinement to choose. All subsequent constructions use the resulting table and its quotient fields.

\begin{proof}[Proof of Theorems~\ref{thm:A} and~\ref{thm:B}]\leavevmode
\begin{enumerate}[label=\emph{Step \arabic*:},ref=\arabic*,leftmargin=4.5em,itemsep=0.4em,beginpenalty=10000]
\item\label{step:deduction-choice} \emph{Choose the construction.} Fix the group, its realization, probability Haar measure and trace inner product, and fix~$\mathsf b>0$ as in Theorem~\ref{thm:B}. Nonabelianness gives~$\dimG>0$. Since~$G\subset U(n_G)$, each~$\operatorname{ad}X$ is skew for~\eqref{eq:gauge-metric}, so
\begin{equation}\label{eq:deduction-casimir}
\casimir=\dimG^{-1}\sum_\alpha\|\operatorname{ad}e_\alpha^{\g}\|_{\rm HS}^2>0\,.
\end{equation}
The sum is positive because~$[\g,\g]\ne0$; hence~$\slope>0$. By Proposition~\ref{prop:U1} there are constants and a construction, chosen in the order of Section~\ref{sec:constants}, including the additional lattice estimates required in Propositions~\ref{prop:U3}, \ref{prop:U4}, \ref{prop:P6}, \ref{prop:P10}, \ref{prop:P12} and~\ref{prop:P13}; fix them.

\item\label{step:deduction-trajectories} \emph{Choose the trajectories.} Keep the level data, path actions, admissible classes and specified recursion of Proposition~\ref{prop:U1}. Proposition~\ref{prop:U2} gives represented trajectories ending at~$\bar\beta$ at level zero, their bare couplings~$\beta_{N,p}$, and the tuning family within this recursion. Its uniform running law gives the slope in B1 of Theorem~\ref{thm:B}.

\item\label{step:deduction-family} \emph{Fix the Wilson family.} Use the local sections and normalized marked transfers of Proposition~\ref{prop:U1}, and set
\begin{equation}\label{eq:deduction-family}
\mathcal T=\bigsqcup_{d\ge0}V_d,\qquad I_{N,p}((d,\mathbf c),y;U)=O^{(d)}_{N,p}(\mathbf c,y;U)\,.
\end{equation}
The insertions are polynomials in the links on the fixed degree-$d$ stencil, with fixed linear renormalization at each~$N,p$. They are continuous and gauge invariant, proving Reg2--Reg4. At~$m=N$, Proposition~\ref{prop:U2} gives Reg1 with~$b_*=\slope>0$, since~$s_N=\log(1/a_N)$. This proves regularity and the remaining lattice assertions B1 of Theorem~\ref{thm:B}. The family, reference levels and token stencils are now fixed.

\item\label{step:deduction-uv} \emph{Take the fixed-volume limits.} Apply Propositions~\ref{prop:U3} and~\ref{prop:U4} to this family. For every~$p,n,v$, obtain the kernel~$f_{p,v}$, smooth on~$\NoHit{L_p}{n}$, with all its fixed finite difference, interpolation and smear limits along the full sequence~$N\to\infty$. This is B2 of Theorem~\ref{thm:B}; its thresholds may depend on~$p$.

\item\label{step:deduction-volume} \emph{Choose one volume refinement.} Let~$\mathsf p:\N\to\N$ be strictly increasing. Proposition~\ref{prop:P5} supplies a strictly increasing~$\mathsf q$, $\mathsf v=\mathsf p\circ\mathsf q$, and one table~$S$, with the field calculus and repeated-word estimate. By Proposition~\ref{prop:P5}(a) the family converges to~$S$ along~$\mathsf v$ in the sense of Definition~\ref{def:ordered-convergence}. Its inner limits are the full limits~$N\to\infty$ of Proposition~\ref{prop:U4} at every admissible volume, with the fixed-volume kernels identified by Proposition~\ref{prop:P5}(b). This gives A-P1 in Theorems~\ref{thm:A} and~\ref{thm:B}, retaining~$\mathsf v$ and~$S$. Since~$L_{\mathsf v(k)}\to\infty$, every fixed compact test support eventually lies in the open fundamental box.

\item\label{step:deduction-stress} \emph{Normalize the stress tensor.} Proposition~\ref{prop:P6} applies with its lattice comparison estimates for the fixed family. It supplies stable Maxwell frames, the conserved line, a nonzero current~$J$ and the raw shell maps. Proposition~\ref{prop:P7} supplies~$c_\star$, the continuous~$O(4)$ actions and the normalized tensor~$T$, with both Ward identities and the joint comparison for all derivatives, arbitrary stress slots, transported derivatives and rotation generators. The scaling sequence of Proposition~\ref{prop:P7}(e) is stable, so its ordinary frames are eventually exactly compatible with conjugation, $H_4$, units and inclusions. Together with Proposition~\ref{prop:P5}, including $\dim_{\C}F_d=\dim_{\C}\maxwellspace{d}$, these give the mixed assertions B3 of Theorem~\ref{thm:B}. Refining comparison radii leaves the volume refinement fixed.

\item\label{step:deduction-phase} \emph{Choose the phase.} Proposition~\ref{prop:P8} gives the compression algebra, its compact metric spectrum~$\phases$, the full-support probability measure~$\nu$ and the phase tables. Each~$S^\xi$ satisfies E1--E7, has null spaces~$\nullspace{d}(S)$ and thus descends to the same~$F_d$, and has the same inclusions, derivatives, conjugation, $O(4)$ representations and stress laws. With the definition of~$\Theta$ in Step~\ref{step:deduction-trace}, this proves B3 of Theorem~\ref{thm:B}. Choose~$\xi_*$ from the full-measure set with one-dimensional OS ground space. This choice imposes no trace condition and is retained below.

\item\label{step:deduction-reconstruction} \emph{Reconstruct the theories.} Apply Proposition~\ref{prop:P9} simultaneously to every~$\xi\in\phases$. Retain its theories and continuations, onto OS comparison maps, unique ground spaces, scalar bounded weak commutants, continued stress tensors, continuity and joint-test averages, and phase-uniform low-energy and heat bounds. These give B4 of Theorem~\ref{thm:B}. For the selected hierarchy use~$\mathfrak W=\mathfrak W^{\xi_*}$ and~$\mathsf F_n=\mathsf F_n^{\xi_*}$. The family and its extraction remain fixed.

\item\label{step:deduction-trace} \emph{Center the trace in the selected phase.} With the tensor and phase just chosen, put
\begin{equation}\label{eq:deduction-trace}
\tau=\sum_{a=1}^4T_{aa},\qquad e_{\rm tr}=S^{\xi_*}(\tau),\qquad \Theta=\tau-e_{\rm tr}\unitfield{4}\,.
\end{equation}
Proposition~\ref{prop:P10} gives its nonzero scalar field, centered operator identity, full-radius asymptotic and unique positive spectral measure, including both high-energy test limits. Proposition~\ref{prop:P11} applies to this same trace and selected theory and gives all three non-freeness conclusions. These are B5 of Theorem~\ref{thm:B}.

\item\label{step:deduction-products} \emph{Choose the products and subtractions.} Proposition~\ref{prop:P12} gives one prescription---channel sets~$\channels{q}$, real fields~$O_{q,\lambda}$, coefficients~$c_{q,\lambda}$ and the associated field projections (Section~\ref{sec:partII-J}.1)---chosen before every spectator, state and truncation, with all its separated, integrated and nested estimates. Proposition~\ref{prop:P13}, with the same projections, gives real stratum bases, scales~$\delta_q^{\rm ct}$ and contact distributions. Its proof uses its own lattice estimates, independently of the bounds of Proposition~\ref{prop:P12}. These give B6 of Theorem~\ref{thm:B}.

\item\label{step:deduction-hierarchy} \emph{Identify the hierarchy.} Take~$F_d=V_d/\nullspace{d}(S)$ and the injective response maps~$\iota_d:F_d\to\mathcal L$ of Proposition~\ref{prop:P5}, and put~$\pi_{\rm quo}(d,\mathbf c)=(d,[\mathbf c])$. Changing a representative in slot~$i$ by~$\mathbf n\in\nullspace{d_i}(S)$ changes~$S_n(v;x)$ by a correlation with~$\mathbf n$ in that slot. By multilinearity, permutation symmetry E1 and the definition of~$\nullspace{d_i}(S)$, this correlation vanishes. The same argument applies in every phase by Step~\ref{step:deduction-phase}. Hence the quotient table satisfies
\begin{equation}\label{eq:deduction-quotient}
S_n^{\rm mix}(\pi_{\rm quo}(v_1),\ldots,\pi_{\rm quo}(v_n);x)=S_n(v;x)\qquad(x\in\Conf{n})\,.
\end{equation}
Use the phase data of Proposition~\ref{prop:P8}, the selected theory and continuation of Proposition~\ref{prop:P9}, the coefficients of Proposition~\ref{prop:P12} and the trace of Proposition~\ref{prop:P10}. These supply H1--H8. On collision diagonals put~$S_n^\zeta=0$. Theorem~\ref{thm:A} evaluates these functions on~$\Conf{n}$: A-P1 through separated tests, R9--R12 and R14 by definition, and~$\tracetwopoint$ at~$r>0$. In O2, the triangle inequality gives~$|x_i-y_j|\ge\delta_{\rm sp}-M_{\rm geo}r>0$ for~$0<r<\delta_{\rm sp}/(M_{\rm geo}+1)$; decrease~$r_1$ accordingly. Internal cluster separation and spectator clearance then put the whole configuration in~$\Conf{n}$. The clause check below proves A-P1--A-P4 for this hierarchy. All assertions of Theorem~\ref{thm:B} now hold. Setting~$\mathsf b=1$ and retaining A-P1--A-P4 proves Theorem~\ref{thm:A}.\qedhere
\end{enumerate}
\end{proof}

\begin{samepage}
\paragraph{The clauses of Theorem~\ref{thm:A}.}
We check the data and each assertion in the order of their definitions.

\emph{Quantifiers and regularity.} Steps~\ref{step:deduction-choice}--\ref{step:deduction-family}, using Propositions~\ref{prop:U1} and~\ref{prop:U2}, fix the strictly increasing positive integer sequences, probability Haar measure and regular family before~$\mathsf p$. Indeed, $D_N=\prod_{i<N}\ell_i$ is strictly increasing because~$\ell_i\ge3$, and~$L_p$ is strictly increasing by Proposition~\ref{prop:U1}. Reg1 has the positive slope~$\slope$, uniformly in~$p$; Reg2 is continuity; Reg3 is gauge invariance; Reg4 uses the fixed finite stencil. Proposition~\ref{prop:P5} then gives one increasing~$\mathsf q$ for each increasing~$\mathsf p$.

\end{samepage}

\emph{H1--H8.} Proposition~\ref{prop:P9} supplies H1 and H6, the selected relativistic datum and continuation. Proposition~\ref{prop:P5} supplies the finite-dimensional normed spaces and response maps H2; their quotient projections supply H3. Proposition~\ref{prop:P8} supplies H4 and, together with Proposition~\ref{prop:P5}, the all-arity functions H5. Propositions~\ref{prop:P12} and~\ref{prop:P10} supply H7 and H8, respectively.

\emph{A-P1.} For a separated test~$f$, set~$C_p^\infty(v;f)=\int f(x)f_{p,v}(x)\,dx$ at every admissible volume, and assign any value at inadmissible volumes. Proposition~\ref{prop:U4} gives the full inner limit there. Eventual admissibility follows from~$L_{\mathsf v(k)}\to\infty$. Proposition~\ref{prop:P5}(a) and~\eqref{eq:deduction-quotient} give the outer limit to~$\int f(x)S_n^{\rm mix}(\pi_{\rm quo}(v_1),\ldots,\pi_{\rm quo}(v_n);x)\,dx$, on the one refinement for every request.

\emph{A-P2, R1.} Proposition~\ref{prop:P9}(a)--(b) gives W0a--c: the translation spectral representation, forward-cone support and unique invariant vacuum line. It gives W1a--e: normalized vacuum, dense domain equal to the field-word span, continuous Schwartz smearing on each domain vector and the adjoint relation. It also gives W2a--d, the label representation law, domain and vacuum invariance and field covariance, and W3, bosonic spacelike locality.

\emph{R2--R6.} Faithfulness in Proposition~\ref{prop:P9}(a) gives R2. The injective response maps of Proposition~\ref{prop:P5}(d) give R3. By E1 in Proposition~\ref{prop:P5}(c), $S_1(\mathbf u_0(x_0))=S_0=1$, so~$\unitfield{0}\ne0$ and~$\iota_0\unitfield{0}\ne0$ by R3; R2 then gives R4. The directed-union property of Proposition~\ref{prop:P5}(d) gives R5, and choosing a representative in each quotient gives the surjectivity R6.

\emph{R7--R12.} Proposition~\ref{prop:P8} gives normalization and full support of~$\nu$ (R7--R8), pointwise phase continuity (R11), and the phase-average identity (R12). Its Euclidean properties, with Proposition~\ref{prop:P5}(c) for the mixture, give permutation symmetry (R9) and smoothness on every~$\Conf{n}$ (R10).

\emph{R13--R15.} Proposition~\ref{prop:P9}(a) gives holomorphy in ordered complex times at fixed real spatial positions, agreement at ordered Euclidean times, and every product-test boundary limit for each strictly ordered height list. R15 retains the stated zero value for a nonintegrable positive-height integral; no extra integrability assertion is needed.

\emph{A-P3, O1--O2.} Proposition~\ref{prop:P12}(a) gives each coefficient bound with its own threshold (O1), and one coefficient system for every spectator, state and truncation (O2), with coefficient~$1$ in both bounds. The remainder threshold is uniform in positions and state.

\emph{A-P4, T1--T4.} Proposition~\ref{prop:P10} gives selected-vacuum centering (T1) and the full positive-radius limit (T3). Its amplitude is~$\traceamplitude=3\dimG/\pi^4>0$, since~$\dimG>0$ (T2). Proposition~\ref{prop:P11}(a) excludes equality with every finite Wick profile on every interval~$(0,\delta)$ (T4).
